% Generated draft: requires mathematical and source review.
\documentclass[11pt]{article}
\usepackage[margin=25mm]{geometry}
\usepackage{fontspec}
\setmainfont{Latin Modern Roman}
\usepackage{amsmath,amssymb,mathtools,unicode-math}
\setmathfont{Latin Modern Math}
\usepackage{enumitem,needspace,xurl,hyperref}
\setlength{\parindent}{0pt}
\setlength{\parskip}{5pt}
\setlength{\emergencystretch}{3em}
\newcommand{\N}{\mathbb{N}}
\newcommand{\Z}{\mathbb{Z}}
\newcommand{\Q}{\mathbb{Q}}
\newcommand{\R}{\mathbb{R}}
\newcommand{\C}{\mathbb{C}}
\newcommand{\F}{\mathbb{F}}
\newcommand{\A}{\mathbb{A}}
\newcommand{\PP}{\mathbb{P}}
\newcommand{\E}{\mathbb{E}}
\newcommand{\eps}{\varepsilon}
\newcommand{\dd}{\,\mathrm{d}}
\begin{document}
\section*{UnsolvedMath ID: 10900010}
\textbf{Source problem code:} catalogue entry
\subsection{Definitions and mathematical statement}
Let \(\Gamma=\pi_1(M)\) act on a simplicial tree \(T\) without a global fixed point. Fix an edge \(e\) and let \(H=\operatorname{Stab}_\Gamma(e)\). Let \(p:M_H\to M\) be the connected cover corresponding to \(H\). The source states that in this cover there is a unique two-dimensional homology class \(\alpha_e\) separating the two ends associated with the two sides of \(T\setminus e\). Its coefficient, orientation, end and compactness conventions are not provided, so existence and uniqueness of \(\alpha_e\) are retained as source input requiring those hypotheses.

Agol asks whether this class has a representative of minimal Thurston norm that embeds in \(M\). Under the natural projection-compatible reading, seek an embedded representative surface \(S\subset M_H\) such that
\[[S]=\alpha_e,\qquad x(S)=x_{M_H}(\alpha_e),\qquad
p|_S:S\to M\text{ is an embedding}.\]
For a compact oriented surface, possibly disconnected, the usual complexity is
\[x(S)=\sum_{C\in\pi_0(S)}\max\{0,-\chi(C)\},\]
and the norm of the class is the minimum of this complexity over its admissible embedded representatives in the covering space. Thus the norm is minimized in the cover while the extra conclusion concerns its projection to the original manifold; mere existence of an incompressible surface or a splitting does not assert that conclusion. The excerpt does not demand connectedness of a norm-minimizing representative.

The source says only ``three-manifold.'' Closedness, orientability, irreducibility, boundary/relative homology, ordinary versus locally finite homology and compact versus properly embedded representatives must be fixed before applying this norm definition. Edge inversions, edge choice and the identification of the cover's end classes also need conventions. If ``embeds in \(M\)'' is meant differently from the covering projection restricting to an embedding, that compatibility must be specified rather than discarded.
\subsection{Short English statement}
Can the separating homology class in an edge-stabilizer cover be represented by a Thurston-norm minimizer whose covering projection is embedded in the original three-manifold? The source omits hypotheses needed to define that class and norm.
\subsection{Sources}
\begin{itemize}
\item Local UnsolvedMath registry, permanent ID \(10900010\), exact source code \texttt{catalogue entry}.
\item Agol, attribution in Delp, Hoffoss and Manning, \emph{Problems In Groups, Geometry, and Three-Manifolds} (2015), subsection 3.2, PDF page 2: \url{https://arxiv.org/pdf/1512.04620#page=2}; abstract \url{https://arxiv.org/abs/1512.04620}. Registry access 2026-07-29; source-tex extraction recorded as checked against rendered PDF text.
\item The summary names Stallings, Thurston norm, Schoen--Yau and Gabai's sutured techniques without exact citations or a theorem proving the stated minimal-norm conclusion.
\end{itemize}
\end{document}
